% answers_short.tex —— 1993 年·旧科目组·理科 · 简短答案速查
% 由 tools/make_short_answers.py 自 answers.tex 的 \gksolhead/\gkanswer 提取，
% 并与 1993/1993_数学真题_旧科目组理_答案要点.md 逐题比对后生成，勿手改（改源 MD 后重跑）。

\begin{gknotice}
\textbf{使用说明}：本节只给最终答案便于快速核对，完整推导见解析版。答案为编写者独立推导，\textbf{非官方标准答案}。
\end{gknotice}

\gkpart{选择题}
\begin{center}\begin{tabular}{ccccccccccccccccccc}
\toprule
题号 & 1 & 2 & 3 & 4 & 5 & 6 & 7 & 8 & 9 & 10 & 11 & 12 & 13 & 14 & 15 & 16 & 17 & 18 \\
\midrule
答案 & C & B & C & D & C & B & B & A & A & D & A & C & D & A & B & B & B & B \\
\bottomrule
\end{tabular}\end{center}

\gkpart{填空题}
\begin{itemize}[leftmargin=2.6em]
  \item[\textbf{19.}] \(2\)
  \item[\textbf{20.}] \(10\sqrt3\approx17.3\)
  \item[\textbf{21.}] \(4186\)
  \item[\textbf{22.}] \(1760\) 元
  \item[\textbf{23.}] \(1\)
  \item[\textbf{24.}] \(\dfrac{1}{a_1d}\)
\end{itemize}

\gkpart{解答题}
\begin{itemize}[leftmargin=2.6em]
  \item[\textbf{25.}] \((0,1)\cup(4,5)\)
  \item[\textbf{26.}] (1) \(A_1C_1\parallel l\)；(2) \(\dfrac{13}{5}\)
  \item[\textbf{27.}] \(\dfrac{x^2}{15/4}+\dfrac{y^2}{3}=1\)
  \item[\textbf{28.}] （条件结论）按题面无解；若分母为 \(1+\overline z^4\)，则 \(\theta=\dfrac{\pi}{12}\) 或 \(\dfrac{7\pi}{12}\)
  \item[\textbf{29.}] (1) 见解析（由 \(p(2)>0\)、\(p(-2)>0\) 及 \(|b|<4\) 推得）；(2) 见解析（由两根分布与 \(|b|<4\) 排除外侧区间）
\end{itemize}
